% Generated draft: requires mathematical and source review.
\documentclass[11pt]{article}
\usepackage[margin=25mm]{geometry}
\usepackage{fontspec}
\setmainfont{Latin Modern Roman}
\usepackage{amsmath,amssymb,mathtools,unicode-math}
\setmathfont{Latin Modern Math}
\usepackage{enumitem,needspace,xurl,hyperref}
\setlength{\parindent}{0pt}
\setlength{\parskip}{5pt}
\setlength{\emergencystretch}{3em}
\newcommand{\N}{\mathbb{N}}
\newcommand{\Z}{\mathbb{Z}}
\newcommand{\Q}{\mathbb{Q}}
\newcommand{\R}{\mathbb{R}}
\newcommand{\C}{\mathbb{C}}
\newcommand{\F}{\mathbb{F}}
\newcommand{\A}{\mathbb{A}}
\newcommand{\PP}{\mathbb{P}}
\newcommand{\E}{\mathbb{E}}
\newcommand{\eps}{\varepsilon}
\newcommand{\dd}{\,\mathrm{d}}
\begin{document}
\section*{UnsolvedMath ID: 701201}
\textbf{Source problem code:} catalogue entry
\subsection{Definitions and mathematical statement}
Let \(F_n\) be the free group with specified free basis \(X=\{x_1,\ldots,x_n\}\). For a reduced word \(a_1\cdots a_m\), where each \(a_i\in X\cup X^{-1}\), its reversal is \(a_m\cdots a_1\); reversal does not invert the letters. An element of \(F_n\) is an \(X\)-palindrome if its unique reduced representative equals its reversal. Write \(\mathcal P_X\) for the set of these elements, including the identity represented by the empty word. Define the palindromic length by
\[
\operatorname{pl}_X(w)=\min\bigl\{k\in\mathbb Z_{\geq 0}:w=p_1\cdots p_k,\quad p_i\in\mathcal P_X\bigr\},
\]
where the empty product is the identity. This minimum exists because each basis letter and its inverse is a palindrome.

\textbf{Question.} Given an element \(w\in F_n\) by a finite word over \(X\cup X^{-1}\), is there an algorithm that always terminates and returns \(\operatorname{pl}_X(w)\)? Equivalently, the input may be required to be reduced.

\textbf{Unresolved scope qualification.} The source gives a finite displayed basis but does not explicitly specify the range of \(n\), or whether the rank is fixed in advance. Under the usual positive finite rank interpretation, the fixed-rank formulation asks whether, for every \(n\geq 1\), there exists such an algorithm for all words in \(F_n\). The uniform formulation asks for a single algorithm taking both \(n\geq 1\) and a word in \(F_n\) as input. These alternatives are recorded without selecting an intended interpretation.

The supplied registry labels the problem open. Its background also reports an affirmative answer attributed to unverified AI solution claims; those claims are not adopted here as established mathematical results.
\subsection{Short English statement}
Relative to a specified free basis, a palindrome is an element whose reduced word reads identically in either direction. The question asks for an algorithm computing the smallest number of palindromes whose product is a given element. The source leaves uniformity in the rank unspecified and reports unverified affirmative solution claims alongside its open label.
\subsection{Sources}
\begin{itemize}
\item Supplied local UnsolvedMath registry record, permanent ID \(701201\), exact source code \texttt{catalogue entry}, title \emph{Kourovka Notebook Problem 16.9}. Proposed in 2006 by V. G. Bardakov, V. A. Tolstykh, and V. E. Shpilrain.
\item Attributed source: \emph{The Kourovka Notebook (Unsolved Problems in Group Theory)}, 16th issue (2006), Problem 16.9; also the supplied 21st issue edition (2026), arXiv:1401.0300v48. The registry describes the collection as edited by E. I. Khukhro and V. D. Mazurov. These source attributions have not been independently checked.
\item The supplied background labels the following as unverified AI solution claims asserting an affirmative answer: Achyuth Jayadevan, \emph{Preprint of 11 September 2026}, \url{https://zenodo.org/records/22722079}; J. Liao, \emph{Preprint of 15 September 2026}, \url{https://arxiv.org/abs/2609.17027}; M. Kinyon and J. Urban, \emph{Preprint of 17 September 2026}, \url{https://kourovkanotebookorg.wordpress.com/wp-content/uploads/2026/09/kourovka-mathematics.pdf}.
\item Supplied status review dated \texttt{2026-10-16}: canonical status \texttt{open}, evidence category \texttt{unrefereed\_claim}. Its notes attribute a \# marker for an unverified AI solution claim to the October 2026 version of the 21st issue, arXiv:1401.0300v48. This is supplied dated curation, not independently verified source or literature assessment. The status-review notes are truncated in the supplied record; the full list of claimed preprints above comes from its background field.
\end{itemize}
\end{document}
